% =============================================================================
% Capítulo 2: Fuentes de Información (nombre literal del enunciado)
% Fuente: docs/fuentes_de_informacion.md (aprobada por el grupo el 05-11-2026).
% =============================================================================
\section{Fuentes de Información}
\label{sec:fuentes}

El enunciado permite usar datasets públicos disponibles en la web. Para cada
línea de negocio se eligió una fuente pública con tres criterios: datos lo más
cercanos posible a una transacción individual, el período más reciente
disponible y una licencia que permita su uso académico.

\subsection{Enfoque: Transacciones Públicas y Clientes Sintéticos}

Los datos públicos de transporte se publican anonimizados, por lo que ningún
dataset identifica a un mismo cliente en distintas líneas de negocio. Como el
caso exige una vista única del cliente, el proyecto combina dos tipos de
datos:

\begin{itemize}[leftmargin=*]
  \item \textbf{Transacciones públicas:} los viajes, pasajes, envíos y
        pedidos provienen de las fuentes de la Tabla~\ref{tab:fuentes}.
  \item \textbf{Clientes sintéticos:} las identidades de los clientes (RUT
        ficticio, género, fecha de nacimiento, tipo de cliente y comuna de
        residencia) se generan de forma sintética. Los clientes residen en
        Chile, y su distribución por comuna, sexo y edad se calibra con el
        Censo de Población y Vivienda 2024 del INE~\cite{ine_censo2024}. Cada
        transacción pública se asigna a uno de estos clientes, de modo que un
        mismo cliente puede comprar en varias líneas.
\end{itemize}

Todo dato sintético se identifica como tal en el código del ETL, en la base de
datos y en este informe.

\subsection{Fuentes por Línea de Negocio}

\begin{table}[H]
  \centering
  \caption{Fuente pública de cada línea de negocio}
  \label{tab:fuentes}
  \small
  \begin{tabularx}{\linewidth}{@{}>{\raggedright\arraybackslash}p{2.4cm}
                                  >{\raggedright\arraybackslash}X
                                  >{\raggedright\arraybackslash}p{3.2cm}
                                  >{\raggedright\arraybackslash}p{1.9cm}
                                  >{\raggedright\arraybackslash}p{2.1cm}@{}}
    \toprule
    \textbf{Línea} & \textbf{Fuente (publicador, país)} & \textbf{Detalle}
      & \textbf{Período} & \textbf{Licencia} \\
    \midrule
    Buses InterUrbanos & Monitriip «Bilhetes de Passagem», Agência Nacional de
      Transportes Terrestres (ANTT), Brasil~\cite{antt_monitriip}
      & Pasajes de buses interestatales por mes, origen, destino y tipo de
        servicio & 2019--2025, mensual & Creative Commons Atribución \\
    Trenes & «UK Train Rides», Maven Analytics, Reino Unido~\cite{maven_uktrain}
      & Un registro por pasaje; datos ficticios según el proveedor
      & Ene.--abr. 2024 & Dominio público \\
    Radio Taxis & TLC Trip Record Data, NYC Taxi \& Limousine Commission,
      EE.UU.~\cite{nyc_tlc} & Un registro por viaje & 2009--2026, mensual
      & Términos de uso de NYC.gov \\
    Logística y Distribución & «Brazilian E-Commerce Public Dataset», Olist,
      Brasil~\cite{olist} & Un registro por pedido e ítem & 2016--2018
      & CC BY-NC-SA 4.0 \\
    Aerolíneas & Airline Origin and Destination Survey (DB1B), Bureau of
      Transportation Statistics, EE.UU.~\cite{bts_db1b}
      & Muestra de pasajes aéreos & Trimestral, incluye 2023--2025
      & Datos públicos del gobierno de EE.UU. \\
    URapid Delivery & «Delivery Center: Food \& Goods orders», Delivery
      Center, Brasil~\cite{gironews_delivery,delivery_center}
      & Un registro por pedido & Ene.--abr. 2021 & No publicada \\
    \bottomrule
  \end{tabularx}
  \nota{Elaboración propia a partir de las páginas de cada fuente, consultadas
  en noviembre de 2026. Para logística se mantiene como alternativa el dataset
  DataCo Smart Supply Chain, con licencia CC BY 4.0~\cite{dataco}.}
\end{table}

La compañía del caso es internacional. Por eso, que las fuentes provengan de
distintos países es coherente con la dimensión \bd{dim_zona}, que registra
país y continente.

\subsection{Uso de cada Fuente en el Modelo}

La Tabla~\ref{tab:uso_fuentes} indica qué datos aporta cada fuente y a qué
parte del modelo estrella llegan.

\begin{table}[H]
  \centering
  \caption{Datos de cada fuente y su destino en el modelo}
  \label{tab:uso_fuentes}
  \small
  \begin{tabularx}{\linewidth}{@{}>{\raggedright\arraybackslash}p{2.4cm}
                                  >{\raggedright\arraybackslash}X
                                  >{\raggedright\arraybackslash}X@{}}
    \toprule
    \textbf{Línea} & \textbf{Datos que aporta} & \textbf{Destino en el modelo} \\
    \midrule
    Buses InterUrbanos & Origen y destino, tipo de servicio, cantidad de
      pasajes, valor medio, gratuidades
      & \bd{dim_zona} (origen y destino), \bd{dim_servicio},
        \bd{cantidad}, montos \\
    Trenes & Estaciones, fecha y hora, precio, tipo de compra (en línea o en
      estación), medio de pago, clase
      & \bd{dim_zona}, \bd{dim_tiempo}, montos, \bd{dim_canal_venta},
        \bd{dim_medio_pago}, \bd{dim_servicio} \\
    Radio Taxis & Zonas de origen y destino, horas de inicio y fin, distancia,
      tarifas, tipo de pago
      & \bd{dim_zona}, \bd{dim_tiempo}, \bd{duracion_minutos},
        \bd{distancia_km}, montos, \bd{dim_medio_pago} \\
    Logística y Distribución & Ciudad del vendedor y del cliente, precio,
      costo de envío, peso del producto, tipo de pago
      & \bd{dim_zona}, montos, \bd{peso_kg}, \bd{dim_medio_pago} \\
    Aerolíneas & Aeropuertos de origen y destino, tarifa, pasajeros,
      trimestre & \bd{dim_zona}, montos, \bd{cantidad}, \bd{dim_tiempo} \\
    URapid Delivery & Canal, medio de pago, monto del pedido, tarifa de envío
      & \bd{dim_canal_venta}, \bd{dim_medio_pago}, montos \\
    \bottomrule
  \end{tabularx}
\end{table}

\subsection{Fuentes de Apoyo}

Además de las fuentes de transacciones, el proyecto usa las fuentes de apoyo de
la Tabla~\ref{tab:fuentes_apoyo}. La dimensión \bd{dim_tiempo} no requiere
fuente externa: se genera con un script SQL.

\begin{table}[H]
  \centering
  \caption{Fuentes de apoyo y su uso en el proyecto}
  \label{tab:fuentes_apoyo}
  \small
  \begin{tabularx}{\linewidth}{@{}>{\raggedright\arraybackslash}p{5.2cm}
                                  >{\raggedright\arraybackslash}X
                                  >{\raggedright\arraybackslash}p{2.4cm}@{}}
    \toprule
    \textbf{Fuente} & \textbf{Uso} & \textbf{Se carga al DWH} \\
    \midrule
    Censo de Población y Vivienda 2024, INE~\cite{ine_censo2024}
      & Calibrar los clientes sintéticos y poblar sus comunas y regiones de
        residencia & Sí, en \bd{dim_cliente} \\
    Tablas de zonas de cada fuente
      & Zonas de taxi de Nueva York, puntos de origen y destino de ANTT,
        estaciones del Reino Unido y aeropuertos de EE.UU.
      & Sí, en \bd{dim_zona} \\
    Tipos de cambio del Banco Central de Chile~\cite{bcentral}
      & Expresar en pesos chilenos (CLP) los montos en dólares, libras y
        reales & Sí, en los montos \\
    Encuesta de transporte interurbano de pasajeros, INE~\cite{ine_etranp};
    memoria anual de EFE~\cite{efe_memoria}; estadísticas de la
    JAC~\cite{jac}
      & Contexto chileno para el análisis y las recomendaciones & No \\
    \bottomrule
  \end{tabularx}
\end{table}

\subsection{Tratamientos Previos a la Carga}

Antes de cargar los datos en el Datawarehouse, el ETL aplica los tratamientos
siguientes. Las limitaciones que estos tratamientos implican se detallan en la
Sección~\ref{sec:limitaciones}.

\begin{enumerate}[leftmargin=*]
  \item \textbf{Clientes sintéticos:} generación de los clientes residentes en
        Chile y asignación de cada transacción a un cliente.
  \item \textbf{Desagregación:} las fuentes que publican totales (ANTT por mes
        y ruta, DB1B por trimestre) se convierten en transacciones
        individuales sintéticas que reproducen esos totales.
  \item \textbf{Período de análisis:} el objetivo pide los últimos tres años,
        por lo que se trabaja con 2023 a 2025. Las fuentes de otros años
        (Olist, Delivery Center) o que cubren solo una parte del período
        (UK Train Rides) trasladan sus fechas a ese rango, conservando su
        estacionalidad.
  \item \textbf{Unidades y moneda:} millas a kilómetros, gramos a kilogramos y
        montos a CLP.
  \item \textbf{Muestreo:} de las fuentes con millones de registros, como los
        viajes de taxi de Nueva York, se carga una muestra.
\end{enumerate}

\subsection{Estado de Verificación de las Fuentes}

Las fuentes se revisaron en noviembre de 2026. Los datos de UK Train Rides,
NYC TLC y DataCo se verificaron en sus páginas oficiales. De ANTT, Olist y BTS
solo se vieron extractos en resultados de búsqueda (el portal de ANTT bloquea
el acceso automatizado y la página de BTS respondió con error), y la página
original de Delivery Center no se encontró: se revisó una copia publicada en
GitHub. Antes de la carga del Avance N°2 se
confirmarán en el navegador las columnas, la granularidad y la licencia de
estas fuentes, además de la serie de tipos de cambio que se usará.
